\documentclass[11pt,a4paper]{article}
\usepackage[utf8]{inputenc}
\usepackage[T1]{fontenc}
\usepackage{lmodern}
\usepackage{textcomp}
\usepackage{amsmath}
\usepackage{array}
\usepackage[hmargin=1.6cm,vmargin=1.2cm]{geometry}
\usepackage{enumitem}
\pagestyle{empty}
\setlength{\parindent}{0pt}
\setlength{\parskip}{0pt}
\begin{document}
\raggedright
{\Large\bfseries Tarea de entornos digitales}
\vspace{10pt}
\begin{enumerate}[label=\textbf{\arabic*.},leftmargin=*,itemsep=6pt,topsep=0pt]
\item \begin{minipage}[t]{\linewidth}
Se dañó el computador de Luis, pero su trabajo estaba guardado en la nube. ¿Qué pasa con el trabajo?

\vspace{4pt}
\begin{tabular}{@{}>{\raggedright\arraybackslash}p{0.47\linewidth}>{\raggedright\arraybackslash}p{0.47\linewidth}@{}}
a)~Solo lo recupera si también lo guardó en un disco duro externo & b)~Lo puede abrir desde otro dispositivo, porque está en servidores a los que accede por internet \\
c)~Solo lo abre desde el mismo computador cuando lo reparen & d)~Se pierde, porque la nube está dentro del computador \\
\end{tabular}
\end{minipage}

\item \begin{minipage}[t]{\linewidth}
La información de internet se actualiza constantemente, lo que permite disponer de contenidos recientes y en permanente evolución. ¿Qué característica de internet es?

\vspace{4pt}
\begin{tabular}{@{}>{\raggedright\arraybackslash}p{0.47\linewidth}>{\raggedright\arraybackslash}p{0.47\linewidth}@{}}
a)~Dinámico & b)~Global \\
c)~Interconectado & d)~Basado en protocolos \\
\end{tabular}
\end{minipage}

\item \begin{minipage}[t]{\linewidth}
Un video con información falsa sobre una vacuna se comparte miles de veces en pocas horas. ¿Qué desventajas de las redes sociales se ven?

\vspace{4pt}
\begin{tabular}{@{}>{\raggedright\arraybackslash}p{0.47\linewidth}>{\raggedright\arraybackslash}p{0.47\linewidth}@{}}
a)~Acortan distancias y son masivas & b)~Las noticias falsas y la viralización de información \\
c)~Permiten compartir información y son canal de aprendizaje & d)~El grooming y el abuso en el uso de las redes \\
\end{tabular}
\end{minipage}

\item \begin{minipage}[t]{\linewidth}
Un adulto engaña a un menor por internet para ganarse su confianza con fines de abuso. ¿Cómo se llama esta desventaja de las redes sociales?

\vspace{4pt}
\begin{tabular}{@{}>{\raggedright\arraybackslash}p{0.47\linewidth}>{\raggedright\arraybackslash}p{0.47\linewidth}@{}}
a)~Noticias falsas & b)~Acceso indiscriminado a contenidos \\
c)~Viralización de información & d)~Grooming \\
\end{tabular}
\end{minipage}

\item \begin{minipage}[t]{\linewidth}
CNT, Claro, Netlife y Xtrim aparecen en el aula como ejemplo de uno de los elementos que se necesitan para conectarse a internet. ¿Cuál?

\vspace{4pt}
\begin{tabular}{@{}>{\raggedright\arraybackslash}p{0.47\linewidth}>{\raggedright\arraybackslash}p{0.47\linewidth}@{}}
a)~Red social horizontal & b)~Módem o router \\
c)~Navegador & d)~Proveedor de acceso a internet \\
\end{tabular}
\end{minipage}

\item \begin{minipage}[t]{\linewidth}
Escribes «contaminación del agua» en la caja del buscador. Según el aula, esas palabras clave en conjunto determinan

\vspace{4pt}
\begin{tabular}{@{}>{\raggedright\arraybackslash}p{0.47\linewidth}>{\raggedright\arraybackslash}p{0.47\linewidth}@{}}
a)~un dominio & b)~un protocolo estandarizado \\
c)~un criterio de búsqueda & d)~una dirección web \\
\end{tabular}
\end{minipage}

\item \begin{minipage}[t]{\linewidth}
Quieres resultados que no contengan la palabra «venta». ¿Qué usas en la búsqueda avanzada?

\vspace{4pt}
\begin{tabular}{@{}>{\raggedright\arraybackslash}p{0.47\linewidth}>{\raggedright\arraybackslash}p{0.47\linewidth}@{}}
a)~La condición «frase exacta», en la primera parte & b)~La condición «sin las palabras», en la primera parte \\
c)~El filtro Safe Search, en la segunda parte & d)~El filtro de región, en la segunda parte \\
\end{tabular}
\end{minipage}

\item \begin{minipage}[t]{\linewidth}
Necesitas artículos científicos sobre diabetes para un trabajo de la universidad. ¿Qué buscador conviene?

\vspace{4pt}
\begin{tabular}{@{}>{\raggedright\arraybackslash}p{0.47\linewidth}>{\raggedright\arraybackslash}p{0.47\linewidth}@{}}
a)~Google Académico, porque se especializa en documentos académicos y científicos & b)~Ninguno, porque la información académica no está en internet \\
c)~El buscador interno de una red social, porque tiene más publicaciones & d)~El buscador interno de una tienda en línea \\
\end{tabular}
\end{minipage}

\item \begin{minipage}[t]{\linewidth}
Según los ejemplos del aula, ¿qué herramienta de almacenamiento en la nube ofrece 15 GB gratis?

\vspace{4pt}
\begin{tabular}{@{}>{\raggedright\arraybackslash}p{0.47\linewidth}>{\raggedright\arraybackslash}p{0.47\linewidth}@{}}
a)~Google Drive & b)~iCloud \\
c)~OneDrive & d)~Dropbox \\
\end{tabular}
\end{minipage}

\item \begin{minipage}[t]{\linewidth}
¿Qué afirmación es correcta sobre las plataformas de gestión del aprendizaje (LMS) y el Sistema de Gestión Académica (SGA)?

\vspace{4pt}
\begin{tabular}{@{}>{\raggedright\arraybackslash}p{0.47\linewidth}>{\raggedright\arraybackslash}p{0.47\linewidth}@{}}
a)~Moodle es un buscador especializado, como Google Académico & b)~El SGA solo se puede abrir desde una computadora de la universidad \\
c)~Moodle es una plataforma LMS de código abierto y el SGA muestra notas, horarios, asistencia y finanzas & d)~Las LMS solo dan acceso al contenido dentro del campus \\
\end{tabular}
\end{minipage}

\end{enumerate}
\end{document}
